% GENERATED FILE. Edit canonical lessons, then run npm run generate:books.
% canonical-source-hash: fnv1a64:2e115aaf2ff71c25
% canonical-lessons: GU-C335-sadhu, GU-C335-srap, GU-C335-bodh, GU-C335-kahevat, GU-C335-sambhalavvu

\chapter{The Wise Old Man}
\label{ch:gu-the-wise-old-man}
\hlchaptermodality{335}

\begin{chapteropening}
\textbf{By the end of this chapter:} \emph{I can say what lesson an old story teaches.}

Complete the last lesson of chapter 335: I can say what lesson an old story teaches.
\end{chapteropening}

\section[\texorpdfstring{\gu{સાધુ} (sādhu)}{સાધુ (sādhu)}]{\gu{સાધુ} (sādhu) — a holy man}

\label{lesson:GU-C335-sadhu}



\begin{quote}
Before the new one: say the Gujarati for to tempt, then the Gujarati for to wander.
\end{quote}

\subsection*{You'll want to know: \gu{સાધુ}}
\textbf{\gu{સાધુ}} — \emph{sādhu} — \textquotedblleft{}a holy man\textquotedblright{}.

\begin{quote}
\gu{ગામ} \gu{બહાર} \gu{એક} \gu{સાધુ} \gu{રહેતો} \gu{હતો}.
\end{quote}

\begin{quote}
\emph{gām bahār ek sādhu raheto hato.}
\end{quote}

\begin{quote}
A holy man lived outside the village.
\end{quote}

A masculine noun.

\begin{cousinweb}
\textbf{\gu{સાધુ}} (\emph{sādhu}) is a tatsama, Sanskrit \emph{sādhu}, \textquotedblleft{}straight, good\textquotedblright{}, and as a noun \textquotedblleft{}a holy man\textquotedblright{} (Monier-Williams). He gives up the world and wanders, or lives apart.
\end{cousinweb}

\begin{grammarlens}[title={What you've built}]
One more word for telling what happened.
\end{grammarlens}

\subsection*{Your turn}
\begin{itemize}
\raggedright
  \item \emph{Say it:} \emph{sādhu}
  \item \emph{Say it:} \emph{sādhu}, once more
  \item \emph{Say it:} \emph{gām bahār ek sādhu raheto hato.} — A holy man lived outside the village.
  \item \emph{Recall:} say the Gujarati for to tempt, then the Gujarati for to wander, then say \emph{sādhu} again
\end{itemize}

\begin{tcolorbox}[breakable,colback=teal!4,colframe=teal!35!black,title={Before you move on}]
What does \gu{સાધુ} mean? (\textquotedblleft{}A holy man\textquotedblright{}.) Say it once more.
\end{tcolorbox}
\section[\texorpdfstring{\gu{શ્રાપ} (śrāp)}{શ્રાપ (śrāp)}]{\gu{શ્રાપ} (śrāp) — a curse}

\label{lesson:GU-C335-srap}



\begin{quote}
Before the new one: say the Gujarati for to wander, then the Gujarati for a holy man.
\end{quote}

\subsection*{You'll want to know: \gu{શ્રાપ}}
\textbf{\gu{શ્રાપ}} — \emph{śrāp} — \textquotedblleft{}a curse\textquotedblright{}.

\begin{quote}
\gu{સાધુએ} \gu{રાજાને} \gu{શ્રાપ} \gu{આપ્યો}.
\end{quote}

\begin{quote}
\emph{sādhue rājāne śrāp āpyo.}
\end{quote}

\begin{quote}
The holy man put a curse on the king.
\end{quote}

A masculine noun.

\begin{cousinweb}
\textbf{\gu{શ્રાપ}} (\emph{śrāp}) goes back to Sanskrit \emph{śāpa}, \textquotedblleft{}a curse\textquotedblright{} (CDIAL 12393); the Sanskrit word has no \emph{r}, and how the Gujarati form got one is not settled. A curse is given: \textbf{\gu{શ્રાપ} \gu{આપ્યો}}.
\end{cousinweb}

\begin{grammarlens}[title={What you've built}]
One more word for telling what happened.
\end{grammarlens}

\subsection*{Your turn}
\begin{itemize}
\raggedright
  \item \emph{Say it:} \emph{śrāp}
  \item \emph{Say it:} \emph{śrāp}, once more
  \item \emph{Say it:} \emph{sādhue rājāne śrāp āpyo.} — The holy man put a curse on the king.
  \item \emph{Recall:} say the Gujarati for to wander, then the Gujarati for a holy man, then say \emph{śrāp} again
\end{itemize}

\begin{tcolorbox}[breakable,colback=teal!4,colframe=teal!35!black,title={Before you move on}]
What does \gu{શ્રાપ} mean? (\textquotedblleft{}A curse\textquotedblright{}.) Say it once more.
\end{tcolorbox}
\section[\texorpdfstring{\gu{બોધ} (bodh)}{બોધ (bodh)}]{\gu{બોધ} (bodh) — a moral (of a story)}

\label{lesson:GU-C335-bodh}



\begin{quote}
Before the new one: say the Gujarati for a holy man, then the Gujarati for a curse.
\end{quote}

\subsection*{You'll want to know: \gu{બોધ}}
\textbf{\gu{બોધ}} — \emph{bodh} — \textquotedblleft{}a moral (of a story)\textquotedblright{}.

\begin{quote}
\gu{આ} \gu{વાર્તાનો} \gu{બોધ} \gu{શું} \gu{છે}?
\end{quote}

\begin{quote}
\emph{ā vārtāno bodh śũ che?}
\end{quote}

\begin{quote}
What is the moral of this story?
\end{quote}

A masculine noun.

\begin{cousinweb}
\textbf{\gu{બોધ}} (\emph{bodh}) is a tatsama, Sanskrit \emph{bodha}, \textquotedblleft{}knowing, understanding\textquotedblright{} (Monier-Williams): what a story makes you understand at the end.
\end{cousinweb}

\begin{grammarlens}[title={What you've built}]
One more word for telling what happened.
\end{grammarlens}

\subsection*{Your turn}
\begin{itemize}
\raggedright
  \item \emph{Say it:} \emph{bodh}
  \item \emph{Say it:} \emph{bodh}, once more
  \item \emph{Say it:} \emph{ā vārtāno bodh śũ che?} — What is the moral of this story?
  \item \emph{Recall:} say the Gujarati for a holy man, then the Gujarati for a curse, then say \emph{bodh} again
\end{itemize}

\begin{tcolorbox}[breakable,colback=teal!4,colframe=teal!35!black,title={Before you move on}]
What does \gu{બોધ} mean? (\textquotedblleft{}A moral (of a story)\textquotedblright{}.) Say it once more.
\end{tcolorbox}
\section[\texorpdfstring{\gu{કહેવત} (kahevat)}{કહેવત (kahevat)}]{\gu{કહેવત} (kahevat) — a proverb}

\label{lesson:GU-C335-kahevat}



\begin{quote}
Before the new one: say the Gujarati for a curse, then the Gujarati for a moral.
\end{quote}

\subsection*{You'll want to know: \gu{કહેવત}}
\textbf{\gu{કહેવત}} — \emph{kahevat} — \textquotedblleft{}a proverb\textquotedblright{}.

\begin{quote}
\gu{મારા} \gu{દાદાએ} \gu{એક} \gu{કહેવત} \gu{કહી}.
\end{quote}

\begin{quote}
\emph{mārā dādāe ek kahevat kahī.}
\end{quote}

\begin{quote}
My grandfather quoted a proverb.
\end{quote}

A feminine noun.

\begin{cousinweb}
\textbf{\gu{કહેવત}} (\emph{kahevat}) is built on \textbf{\gu{કહેવું}}, to say, from Sanskrit \emph{kathayati}, \textquotedblleft{}converses, describes\textquotedblright{} (CDIAL 2703): a saying everybody knows. It is feminine: \textbf{\gu{એક} \gu{કહેવત} \gu{કહી}}.
\end{cousinweb}

\begin{grammarlens}[title={What you've built}]
One more word for telling what happened.
\end{grammarlens}

\subsection*{Your turn}
\begin{itemize}
\raggedright
  \item \emph{Say it:} \emph{kahevat}
  \item \emph{Say it:} \emph{kahevat}, once more
  \item \emph{Say it:} \emph{mārā dādāe ek kahevat kahī.} — My grandfather quoted a proverb.
  \item \emph{Recall:} say the Gujarati for a curse, then the Gujarati for a moral, then say \emph{kahevat} again
\end{itemize}

\begin{tcolorbox}[breakable,colback=teal!4,colframe=teal!35!black,title={Before you move on}]
What does \gu{કહેવત} mean? (\textquotedblleft{}A proverb\textquotedblright{}.) Say it once more.
\end{tcolorbox}
\section[\texorpdfstring{\gu{સંભળાવવું} (sambhaḷāvvũ)}{સંભળાવવું (sambhaḷāvvũ)}]{\gu{સંભળાવવું} (sambhaḷāvvũ) — to recite; to tell (a story) aloud}

\label{lesson:GU-C335-sambhalavvu}



\begin{quote}
Before the new one: say the Gujarati for a moral, then the Gujarati for a proverb.
\end{quote}

\subsection*{You'll want to know: \gu{સંભળાવવું}}
\textbf{\gu{સંભળાવવું}} — \emph{sambhaḷāvvũ} — \textquotedblleft{}to recite; to tell (a story) aloud\textquotedblright{}.

\begin{quote}
\gu{સાધુએ} \gu{બાળકોને} \gu{એક} \gu{વાર્તા} \gu{સંભળાવી}.
\end{quote}

\begin{quote}
\emph{sādhue bāḷkone ek vārtā sambhḷāvī.}
\end{quote}

\begin{quote}
The holy man told the children a story.
\end{quote}

Past, with \gu{તેણે} (he or she …): \gu{તેણે} \gu{વાર્તા} \gu{સંભળાવી}, \emph{teṇe vārtā sambhḷāvī}.

\begin{cousinweb}
\textbf{\gu{સંભળાવવું}} (\emph{sambhaḷāvvũ}) is the causative of \textbf{\gu{સાંભળવું}}, to hear, from \emph{*saṁbhalati}, \textquotedblleft{}attends to\textquotedblright{}, a form CDIAL reconstructs (CDIAL 12962): to make others hear, to tell a story aloud. The long \emph{ā} of \textbf{\gu{સાંભળવું}} shortens in the longer verb.
\end{cousinweb}

\begin{grammarlens}[title={What you've built}]
One more word for telling what happened.
\end{grammarlens}

\subsection*{Your turn}
\begin{itemize}
\raggedright
  \item \emph{Say it:} \emph{sambhaḷāvvũ}
  \item \emph{Say it:} \emph{sambhaḷāvvũ}, once more
  \item \emph{Say it:} \emph{sādhue bāḷkone ek vārtā sambhḷāvī.} — The holy man told the children a story.
  \item \emph{Recall:} say the Gujarati for a moral, then the Gujarati for a proverb, then say \emph{sambhaḷāvvũ} again
\end{itemize}

\begin{tcolorbox}[breakable,colback=teal!4,colframe=teal!35!black,title={Before you move on}]
What does \gu{સંભળાવવું} mean? (\textquotedblleft{}To recite; to tell (a story) aloud\textquotedblright{}.) Say it once more.
\end{tcolorbox}
